% ── Plantilla UMA · Premios Cátedra de Ciberseguridad ──────────────
\usepackage{graphicx}
\usepackage{xcolor}
\definecolor{umablue}{RGB}{0,59,113}
\definecolor{umacyan}{RGB}{0,174,231}
\definecolor{semgreen}{RGB}{46,125,50}
\definecolor{semyellow}{RGB}{230,150,0}
\definecolor{semred}{RGB}{198,40,40}

\makeatletter
\AtBeginDocument{%
  \hypersetup{colorlinks=true, linkcolor=umablue, urlcolor=umablue, citecolor=umablue}%
  \urlstyle{tt}%
  \g@addto@macro{\UrlBreaks}{\UrlOrds}%
}
\makeatother

% Tipografía de títulos: sans, azul corporativo (sin titlesec)
\makeatletter
\renewcommand\section{\@startsection{section}{1}{\z@}%
  {-12pt \@plus -2pt \@minus -2pt}{5pt}%
  {\Large\bfseries\sffamily\color{umablue}}}
\renewcommand\subsection{\@startsection{subsection}{2}{\z@}%
  {-9pt \@plus -2pt \@minus -2pt}{4pt}%
  {\large\bfseries\sffamily\color{umablue}}}
\renewcommand\subsubsection{\@startsection{subsubsection}{3}{\z@}%
  {-9pt \@plus -2pt \@minus -2pt}{4pt}%
  {\normalsize\bfseries\sffamily\color{umablue!85!black}}}
\makeatother

% Cabecera y pie con logo UMA
\usepackage{fancyhdr}
\AtEndDocument{\label{umaLastPage}}
\setlength{\headheight}{32pt}
\pagestyle{fancy}
\fancyhf{}
\fancyhead[L]{\footnotesize\sffamily\color{black!55}WatchGate · Memoria del proyecto}
\fancyhead[R]{\raisebox{-0.35\height}{\includegraphics[height=0.95cm]{plantilla/logo_uma.png}}}
\fancyfoot[L]{\footnotesize\sffamily\color{black!55}Premios de la Cátedra de Ciberseguridad · Universidad de Málaga · 2026}
\fancyfoot[R]{\footnotesize\sffamily\color{black!55}\thepage\ de \pageref*{umaLastPage}}
\renewcommand{\headrulewidth}{0.6pt}
\renewcommand{\headrule}{\hbox to\headwidth{\color{umacyan}\leaders\hrule height \headrulewidth\hfill}}
\renewcommand{\footrulewidth}{0.3pt}
\renewcommand{\footrule}{\hbox to\headwidth{\color{black!25}\leaders\hrule height \footrulewidth\hfill}}

% Diagramas
\usepackage{tikz}
\usetikzlibrary{positioning,arrows.meta}
\usepackage{float}

% Tablas y espaciado
\usepackage{booktabs}
\usepackage{etoolbox}
\AtBeginEnvironment{longtable}{\small\sffamily}
\setlength{\emergencystretch}{3em}
\AtBeginDocument{\setlength{\parskip}{3pt plus 1pt}}

% Bloques de código algo más compactos
\AtBeginEnvironment{verbatim}{\footnotesize}
